%% \glt (the free translation) and its style hook \GlossTransStyle.
%% The hook is a DECLARATION, applied after \glt's \par, so it must
%%   (a) reach the translation,
%%   (b) leave the gloss tiers above it untouched,
%%   (c) not leak past the end of the example.
%% Every sentinel below appears TWICE with identical spelling, once in the
%% unstyled example and once in the styled one, so the two occurrences are
%% width-comparable: same glyphs, so any width difference is the hook and
%% nothing else.  A size declaration is used rather than \itshape because
%% italic vs upright differs by only ~1pt in Latin Modern, too thin to
%% assert on across three engines; \Large is unmistakable.
\input{_preamble}
\begin{document}
GLTBODY reference line before any example.

%% unstyled reference: the hook is empty by default
\exg. Der GLTOBJ bellte.\\ the GLTGLOSS bark.PST\\
\glt GLTTRANS barked here.

\renewcommand\GlossTransStyle{\Large}

%% identical gloss, styled translation
\exg. Der GLTOBJ bellte.\\ the GLTGLOSS bark.PST\\
\glt GLTTRANS barked here.

GLTBODY reference line after the styled example.
\end{document}
